\section{Fundamentação Teórica}

Considere uma gotícula de raio $r$ feita de um material de densidade $\sigma$.
Se $g$ é a aceleração da gravidade, o módulo da força gravitacional agindo sobre ela é dado por
\begin{equation}
F_1 = mg = \frac{4}{3} \pi r^3 \sigma g \;.
\label{eq: F_gravitational}
\end{equation}

Se a gota está imersa em um meio de densidade $\rho$ e viscosidade $\eta$,
e se move com velocidade escalar $v$, temos também a força de empuxo e a de arrasto, 
dadas respectivamente por
\begin{equation}
F_2 = \frac{4}{3} \pi r^3 \rho g \;,
\label{eq: F_buoyancy}
\end{equation} e
\begin{equation}
F_3 = 6 \pi r \eta v \;.
\label{eq: F_drag}
\end{equation}

Considere que a gota se encontra entre as placas de um capacitor de placas paralelas. 
Se $V$ é a tensão do capacitor, e $d$ a distância entre suas placas, então o campo elétrico 
na região entre elas é constante e tem módulo $E = V/d$. Seja $q$ o módulo da carga na 
gotícula, então o módulo da força elétrica agindo sobre ela é 
\begin{equation}
F_4 = qE = qV/d \;.
\label{eq: F_electric}
\end{equation} 

O sentido de movimento será dado pelo sentido da força elétrica. A força de arrasto terá sentido 
oposto ao movimento. A força gravitacional sempre aponta para baixo, e o empuxo, para cima.

Vamos supor que as gotas atingem rapidamente a velocidade terminal, de modo que 
seu movimento é retilíneo e uniforme. i.e., $\sum_i \mathbf F_i = 0$. 
Tomando o sentido para cima como sendo positivo, temos que quando a gotícula se move para cima, 
a condição de equilíbrio pode ser exprimida por 
\begin{gather}
-F_1 +F_2 -F_3 +F_4 = 0 \notag \\
\implies - {\frac{4}{3} \pi r^3 \sigma g} + {\frac{4}{3} \pi r^3 \rho g} - {6 \pi r \eta v_s} + {qV/d} = 0 \notag \\
\implies v_s = \frac{1}{6 \pi r \eta} \left[ qV/d + \frac{4\pi}{3} r^3 g (\rho - \sigma) \right] \;,
\label{eq: v_subida}
\end{gather}
\noindent onde $v_s$ é a velocidade de subida da gotícula. 
Analogamente, se $v_d$ é a velocidade de descida, temos
\begin{gather}
-F_1 +F_2 + F_3 - F_4 = 0 \notag \\
\implies - {\frac{4}{3} \pi r^3 \sigma g} + {\frac{4}{3} \pi r^3 \rho g} + {6 \pi r \eta v_d} -{qV/d}   \notag \\
\implies v_d = \frac{1}{6 \pi r \eta} \left[ qV/d - \frac{4\pi}{3} r^3 g (\rho - \sigma) \right] = 0\;.
\label{eq: v_descida}
\end{gather}

Subtraíndo a eq. \eqref{eq: v_descida} da eq. \eqref{eq: v_subida}, obtemos
\begin{align}
v_s - v_d &= \frac{2}{6 \pi r \eta} \frac{4 \pi r^3 g}{3} (\rho - \sigma) \notag = \frac{4g}{9\eta} r^2 (\rho - \sigma) \notag \\[2pt]
\implies r^2 &= \frac{9 \eta}{4 g} \frac{v_s - v_d}{\rho - \sigma} \notag \\
\implies r &= \frac{3}{2} \left[ \frac{\eta}{g|\sigma - \rho|} \right]^{1/2} |v_s - v_d|^{1/2} \notag \\
&= A |v_s - v_d|^{1/2}\;.
\label{eq: r}
\end{align}
onde
\begin{equation}
A \equiv \frac{3}{2} \left[ \frac{\eta}{g|\sigma - \rho|} \right]^{1/2} \;.
\end{equation}

Somando as eqs. \eqref{eq: v_subida} e \eqref{eq: v_descida}, obtemos
\begin{align}
v_s + v_d = 2\left( \frac{qV}{6 \pi r \eta d} \right) = \frac{qV}{3 \pi r \eta d} \notag \\
\implies q = 3\pi r \eta d (v_s + v_d)/V \;.
\label{eq: q}
\end{align}

Substituindo a eq. \eqref{eq: r} na \eqref{eq: q}, obtemos
\begin{gather}
q = 3\pi d \eta \frac{A}{V} (v_s + v_d) |v_s - v_d|^{1/2} \;.
\label{eq: q'}
\end{gather}
